% Authority: EGA II front-fr.tex, SHA-256 F199DE0C2C8C49E2BDAEF48D764086952806E77D8066795B97B3595F7EA03198.
\clearpage
\thispagestyle{plain}
\markboth{제II장}{제II장}
\oldpage[II]{5}
\phantomsection
\label{chapter:II-ko}

\begin{center}
{\large 제II장\par}
\vspace{1.5em}
{\Large\bfseries 몇몇 사상 부류에 대한\par}
{\Large\bfseries 초보적 전역 연구\par}
\vspace{1.8em}
{\bfseries 목차\par}
\end{center}

\medskip
\begin{tabular}{@{}r@{\ }p{0.82\textwidth}@{}}
\S~1. & 아핀 사상.\\
\S~2. & 동차 소 스펙트럼.\\
\S~3. & 등급 대수의 층의 동차 스펙트럼.\\
\S~4. & 사영 다발. 풍부한 층.\\
\S~5. & 준아핀 사상; 준사영 사상; 고유 사상; 사영 사상.\\
\S~6. & 정수적 사상과 유한 사상.\\
\S~7. & 값매김 판정법.\\
\S~8. & 블로업 스킴; 사영 원뿔; 사영 폐포.
\end{tabular}

\medskip
이 장에서 다루는 여러 사상 부류는 코호몰로지적 방법을 크게 사용하지
않고 연구한다. 이 방법을 사용하여 그 사상 부류들을 더 깊이 연구하는
일은 제III장에서 이루어질 것이며, 거기서는 제II장의 \S\S~2, 4, 5를
주로 사용할 것이다. \S~8은 처음 읽을 때 생략해도 된다. 이 절은
\S\S~1에서~3까지 전개한 형식 체계에 몇 가지 보충을 더하는데, 그
내용은 이 형식 체계의 쉬운 응용으로 귀착된다. 우리는 이 절의 결과들을
이 장의 다른 결과들만큼 자주 사용하지는 않을 것이다.
